\documentclass[11pt]{article}
\usepackage{manual_docs_style}
\externaldocument{build/candidate_generation}[candidate_generation.pdf]
\externaldocument{build/dataset_manifest_schema}[dataset_manifest_schema.pdf]

\begin{document}

\docTitle{Web Explorers}
\phantomsection\label{docstart:web_model_explorer}

This document describes the local web explorers used to inspect \name{} runs and generated question bundles.

\section*{Model Explorer}

The Web Model Explorer is a local, read-only browser for inspecting  materialized Simulink runs.
The registry is loaded with a paged approach to handle large policy (i.e. the design principle is that we fetch only what is needed).
The first page contains 50 families, and additional pages are fetched on demand (by clicking next page)
Full family details, including intervention and probe runs, are loaded when a family is expanded (i.e. when we are clicking on a family)
Deep links are supported and can open a selected run without requiring the full policy registry to be loaded up front.

Start it from the repository with:

\begin{DocCode}
bash web_model_explorer/start.sh
\end{DocCode}

The launcher starts the Next.js app in \code{web_model_explorer/}.
It uses \code{0.0.0.0:3001} by default.
\code{WEB_MODEL_EXPLORER_HOST}, \code{WEB_MODEL_EXPLORER_PORT}, and \code{WEB_MODEL_EXPLORER_RUNS_DIR_NAME} can override the host, port, and displayed run-artifact folder.

\section*{What It Shows}

The sidebar lists allowed \name{} \termSimulators that exist under \code{models/simulink/}.
\termSimulators with missing required files remain listed but show a warning badge.
The header shows the selected \termSimulator, selected \code{policy_id}, active runs folder name, and a reload button. 
The reload button reloads the visible page of families and any expanded rows.

The main table is organized by \code{family_id} from the resolved run graph.
Each family can be expanded to show its baseline run, intervention runs, and matched time-zero/static-change probe runs.

Table columns include:

\begin{itemize}
\item status from \code{model_record.json};
\item cheap-filter status from \code{cheap_filter_metrics.jsonl};
\item \code{run_id};
\item \code{family_id};
\item \code{source};
\item \code{validation_profile};
\item \code{recipe_hash};
\item exposed parameter and initial-state values in experiment-config order;
\item actions for plotting and expansion.
\end{itemize}

Expanded intervention rows show changed parameter, direction, new value, and intervention time from \code{run_nodes.jsonl} and \code{run_edges.jsonl}.
For production policies with surrogate scoring, the table also shows \code{surrogate_id}, \code{noise_level}, \code{predicted_label}, \code{true_label_confidence}, \code{confidence_margin}, and \code{surrogate_filter_pass}.

\section*{Plot Area}

The plot area renders selected runs.
Plot data can be shown with no noise or with a \termSimulator-provided low or high noise profile when \code{noise_adder.py} exists.

Controls include:

\begin{itemize}
\item run-selection eye buttons in the table;
\item noise profile: \code{none}, \code{low}, or \code{high};
\item deterministic noise seed;
\item observed signal selection.
\end{itemize}

If two runs are selected and noise is enabled, the UI may show per-signal SNR values using the \termSimulator noise-adder utilities.

\section*{Deep Links}
Deep links can preselect a \termSimulator, policy, run, and optional comparator.
When a deep link targets a run, the explorer resolves the run and fetches only the containing family needed to display it, rather than loading the full policy registry.
Supported query parameters include:

\begin{DocCode}
model, policy, run, run2, compare_run, compare, noise_profile, noise_seed
\end{DocCode}

The \code{compare} parameter may be \code{baseline}, \code{time0}, \code{sibling}, or \code{none}.
A plain \code{?run=<RUN_ID>} link opens a single-run view unless an explicit comparator is provided.

Shortcut:

\begin{DocCode}
http://localhost:3001/<run_id>
\end{DocCode}

opens the matching run directly if it can be located.

\section*{Endpoints Used}

\begin{itemize}
\item \code{GET /api/models}: read allowed \termSimulators and validation warnings.
\item \code{GET /api/policies?model=<simulator_id>}: list available resolved plans under \code{plans/}.
\item \code{GET /api/registry?model=<simulator_id>&policy=<POLICY_ID>}: load the first page of family summaries for the selected policy.
\item \code{GET /api/registry?model=<simulator_id>&policy=<POLICY_ID>&page=<N>&page_size=<N>}: load a specific page of family summaries.
\item \code{GET /api/registry?model=<simulator_id>&policy=<POLICY_ID>&family_id=<FAMILY_ID>}: load full details for one family, including baseline, intervention, and probe runs.
\item \code{GET /api/registry?model=<simulator_id>&policy=<POLICY_ID>&run=<RUN_ID>}: resolve a run within a policy and load the containing family for deep-link display.
\item \code{GET /api/registry?model=<simulator_id>&policy=<POLICY_ID>&mode=full}: explicitly load the full registry for debugging or compatibility.
\item \code{GET /api/distribution?model=<simulator_id>}: load \code{experiment_config.json} and generated metadata used for legal domains and sampling details.
\item \code{GET /api/locate-run?run=<RUN_ID>}: resolve a deep-linked run identifier to \termSimulator and policy.
\item \code{GET /api/data?model=<simulator_id>&run=<RUN_ID>&max_rows=<N>&noise_profile=<PROFILE>&noise_seed=<SEED>}: load a run artifact for plotting.
\item \code{GET /api/simulate?model=<simulator_id>&policy=<POLICY_ID>}: read saved simulation status information when present.
\end{itemize}

\section*{Files Read}

\begin{itemize}
\item \code{models/simulink/<simulator_id>/plans/<policy_id>/run_nodes.jsonl}
\item \code{models/simulink/<simulator_id>/plans/<policy_id>/run_edges.jsonl}
\item \code{models/simulink/<simulator_id>/plans/<policy_id>/generation_report.json}
\item \code{models/simulink/<simulator_id>/plans/<policy_id>/validation_report.json}
\item \code{models/simulink/<simulator_id>/<runs>/model_record.json}
\item \code{models/simulink/<simulator_id>/metrics/<policy_id>/cheap_filter_metrics.jsonl}
\item \code{models/simulink/<simulator_id>/metrics/<policy_id>/surrogate_scores/<surrogate_id>.jsonl}
\item \code{models/simulink/<simulator_id>/surrogates/<surrogate_id>/thresholds.json}
\item \code{models/simulink/<simulator_id>/generated/metadata.json}
\item \code{models/simulink/<simulator_id>/description_levels.json}
\item \code{models/simulink/<simulator_id>/experiment_config.json}
\item \code{models/simulink/<simulator_id>/<runs>/<run_id>/data.parquet}
\item \code{tsENV_questions/<simulator_id>/dataframes/<run_id>.parquet}, as fallback for exported bundles
\item \code{models/simulink/<simulator_id>/noise_adder.py}
\end{itemize}


\section*{Web Human Study}

The Web Human Study app checks whether exported questions are feasible for humans.
It reads from \code{tsENV_questions/} and can be started with:

\begin{DocCode}
bash web_human_study/start.sh [--port PORT]
\end{DocCode}

The app visualizes time series with Plotly and shows the multiple-choice labels from \code{questions.json}.
When a user selects an option, the correct label is shown, and the result is saved under \code{human_feedback/<name>.json}.
There is also a button to mark a question as problematic.

Shortcut forms include:

\begin{itemize}
\item \code{http://localhost:5173/<run_id>}: open a sample by \termRunID.
\item \code{http://localhost:5173/<question_id>}: open a question.
\item \code{http://localhost:5173/<question_id>-<run_id>}: open a specific run within a question.
\end{itemize}

\end{document}
